
This work investigated sequential feedback routing for LLM code generation, comparing cascade routing (static analysis before conditional execution) against aggregation baseline (simultaneous multi-source feedback).

\textbf{Validated findings:}
\begin{itemize}
\item 35 dual-sensitive tasks identified from HumanEval (21.3\% of benchmark)
\item Mypy --strict detects errors in 99.6\% of CodeLlama-7B generation attempts on dual-sensitive tasks
\item Mock simulation suggests token efficiency ratio of 0.733 for cascade vs. aggregation
\end{itemize}

\textbf{Unvalidated claims:}
\begin{itemize}
\item Iteration reduction via attention economy (implementation failure prevented testing)
\item Token efficiency with real model inference (mock simulation only)
\end{itemize}

\textbf{Scope limitations:}
\begin{itemize}
\item Results specific to CodeLlama-7B base model, dual-sensitive HumanEval tasks, Python with mypy --strict
\item Detection rate may not generalize to instruction-tuned models or models with stronger type safety training
\end{itemize}

\textbf{Implications:}
The extreme mypy detection rate (99.6\%) suggests that for certain model-task combinations, static analysis serves as an effective pre-filter, catching nearly all type-related errors before execution. This validates the computational efficiency premise of cascade routing: conditional execution gating can skip expensive testing in the majority of iterations.

The attention economy hypothesis---whether sequential presentation reduces LLM cognitive load---remains an open question for future work.
